\documentclass{article}
\usepackage{amsmath}
\usepackage{tikz}
\usetikzlibrary{positioning,calc,decorations.pathreplacing,shapes,fit}
\usepackage{colortbl}
\usepackage{booktabs}
\usepackage{url}
\usepackage{pifont}
\newcommand{\cmark}{\ding{51}}%
\newcommand{\xmark}{\ding{55}}%
\usepackage{array, adjustbox,url}
\usepackage{pifont,marvosym} % wasysym
\usepackage{rotating,subfig}
\usepackage{xspace}

\title{BTI 3021: Networking Project - Sprint 3}

\author{Christian Grothoff}
\date{}

\begin{document}
\maketitle

\section{Introduction}

For this sprint you will write an IP router, building upon the
results from your sprint.

While the driver and skeleton you are given is written in C, you may
again use {\em any} language of your choice for the implementation (as
long as you extend the {\tt Makefile} with adequate build rules).
However, if you choose a different language, be prepared to write
additional boilerplate yourselves.

How an IP router works is expected to be understood from the
networking class. If not, you can find plenty of documentation and
specifications on the Internet.

The basic setup is the same as in the first two sprints.

\subsection{Deliverables}

There will be three main deliverables for the sprint:
\begin{description}
\item[Documentation] You must document the main
  data structure (routing table) and a {\bf comprehensive test plan}.
\item[Implementation] You must implement the IP protocol. Your
  implementation must answer to IP packets, and route them.
  For this, you are to extend the {\tt router.c} template provided
  (or write the entire logic from scratch in another language).
\item[Testing] You must implement and submit your own {\bf test cases}
  by {\em pretending} to be the network driver (see below) and sending
  IP packets or command-line inputs to your program and verifying that it
  outputs the correct frames. Additionally, you should perform and
  {\bf document} {\em interoperability} tests against existing
  implementations (i.e. other notebooks from your team to ensure that
  your IP router implementation integrates correctly with other
  implementations).
\end{description}

All deliverables must be submitted to Git (master branch)
by {\bf January 22nd 2021, 18:00 CET}.


\subsection{Team}

You are expected to work in a team of four students. The suggested
division of work is:
\begin{description}
\item[1 Student] Scrum master, also responsible for the documentation
\item[1 Student] Implementation
\item[2 Students] Testing
\end{description}


\subsection{Grading}

You get points for each of the key deliverables:
\begin{center}
\begin{tabular}{l|r}
Technical documentation (A\&D, test plan) &  3 \\ \hline
Correct implementation                    &  9 \\ \hline
Comprehensive test cases                  &  4 \\ \hline
Quality of English in documentation       &  3 \\ \hline \hline
Total                                     & 19
\end{tabular}
\end{center}

We will specifically also look for the following properties of a router:

\begin{itemize}
\item Correct forwarding? % IP packets flow? Mac updated?
\item Correct address resultion and caching? % ARP cache?
\item Correct IP handling? % IP TTL decremented, ICMP? Checksum?
\item Correct IP fragmentation? % Use eth3 for testing
\end{itemize}

Parial points are awarded.

You can earn 16, 15 and 19 points in the three networking sprints, for
a total of 50 points.  You need {\bf 37 points} to pass the networking
component of BTI 3021.


\section{Detailed interface specification}

Implement {\tt router} which routes IPv4 packets:
\begin{enumerate}
\item Populate your routing table from the network interface configuration
  given on the command-line using the same syntax as with the {\tt arp}
  program.
\item Use the ARP logic to resolve the target MAC address.  You MUST
  simply drop IP packets for destinations where you do not yet have the
  next hop's MAC address, but you MUST then issue the ARP request to
  obtain the next hop or destination's MAC instead (once per dropped IP packet).
\item Make sure to decrement the TTL field and recompute the CRC.
  % add link to logic implementing CRC?
\item Generate ICMP messages for ``no route to network'' (ICMP
    Type 3, Code 0) and ``TTL exceeded'' (ICMP Type 11, Code 0),
\item Support the syntax {\tt IFC[RO]=MTU} where {\tt MTU} is the
  MTU for IFC.  Example: {\tt eth0=1500}.  Implement and test IPv4 fragmentation
  (including {\em do not fragment}-flag support), including sending
  ICMP  (ICMP Type 3, Code 4).
\item Support dynamic updates to the routing table via {\tt stdin}.
  Base your commands on the {\tt ip route} tool.  For example,
  ``route list'' should output the routing table, and
  ``route add 1.2.0.0/16 via 192.168.0.1 dev eth0'' should add
  a route to {\tt 1.2.0.0/16} via the next hop {\tt 192.168.0.1}
  which should be reachable via {\tt eth0}.  Implement at least
  the {\tt route list}, {\tt route add} and {\tt route del} commands.
\end{enumerate}

% You do not need to support VLANs, IP multicast or IP broadcast.


%Configure your router with {\tt eth1} using 192.168.0.1/16.  Configure
%{\tt eth2} using 10.0.0.1/8 and {\tt eth3} using 172.16.0.1/12.
%Connect your notebook as 10.0.0.2 using 10.0.0.1 as the default
%route. Set an MTU of 500 on {\tt eth3}.  Set a default route of
%192.168.0.2 on the router.


\end{document}
